\documentclass[12pt]{article}
\usepackage{amsmath,amssymb,amsfonts,mathrsfs}
\usepackage[margin=1in]{geometry}

\begin{document}

\section*{Chapter 2, Problem 12}

\subsection*{Problem Statement}
A Lagrangian density for classical electrodynamics is
\[
  \mathscr{L} = -\frac{1}{4}\sum_{\mu,\nu} F_{\mu\nu}F_{\mu\nu} + \frac{1}{c}\sum_\mu j_\mu A_\mu,
\]
with
\[
  F_{\mu\nu} = \frac{\partial A_\nu}{\partial x_\mu} - \frac{\partial A_\mu}{\partial x_\nu}.
\]
The four-vectors $\{A_\mu\}$ and $\{j_\mu\}$ are defined by
\[
  \{A_\mu\} = \{A_1, A_2, A_3, A_4\} = \{{\bf A}, i\phi\},
\]
\[
  \{j_\mu\} = \{j_1, j_2, j_3, j_4\} = \{{\bf j}, ic\rho\},
\]
where ${\bf A}$ is the vector potential (${\bf B} = \nabla\times{\bf A}$), $\phi$ is the scalar potential
\[
  \left({\bf E} = -\nabla\phi - \frac{1}{c}\frac{\partial{\bf A}}{\partial t}\right),
\]
${\bf j}$ is the current vector, and $\rho$ is the charge density (see Problem 1--13). Show that this Lagrangian density leads to the familiar Maxwell equations.

Using Lagrange's equations with ${\bf A}$ and $\phi$ as dependent variables, show that this Lagrangian density leads to the familiar Maxwell equations.

\subsection*{Solution}

The independent variables are $x_\mu = (x_1, x_2, x_3, x_4) = (x, y, z, ict)$, and the dependent field variables are the four components $A_\mu$.

The Euler--Lagrange equation for $A_\lambda$ is:
\[
  \frac{\partial\mathscr{L}}{\partial A_\lambda} - \sum_\mu \frac{\partial}{\partial x_\mu}\frac{\partial\mathscr{L}}{\partial(\partial A_\lambda/\partial x_\mu)} = 0.
\]

\subsubsection*{Step 1: Compute $\partial\mathscr{L}/\partial A_\lambda$}

The only term in $\mathscr{L}$ that depends on $A_\lambda$ (not its derivatives) is $\frac{1}{c}j_\mu A_\mu$:
\[
  \frac{\partial\mathscr{L}}{\partial A_\lambda} = \frac{1}{c}j_\lambda.
\]

\subsubsection*{Step 2: Compute $\partial\mathscr{L}/\partial(\partial A_\lambda/\partial x_\mu)$}

Since $F_{\alpha\beta} = \partial A_\beta/\partial x_\alpha - \partial A_\alpha/\partial x_\beta$, the derivative $\partial A_\lambda/\partial x_\mu$ appears in:
\begin{itemize}
  \item $F_{\mu\lambda}$ (from $\partial A_\lambda/\partial x_\mu$), contributing $+1$;
  \item $F_{\lambda\mu}$ (from $-\partial A_\lambda/\partial x_\mu$... wait, $F_{\lambda\mu} = \partial A_\mu/\partial x_\lambda - \partial A_\lambda/\partial x_\mu$), contributing $-1$ to the $-\partial A_\lambda/\partial x_\mu$ term.
\end{itemize}

Working through carefully:
\[
  \frac{\partial}{\partial(\partial A_\lambda/\partial x_\mu)}\left(-\frac{1}{4}F_{\alpha\beta}F_{\alpha\beta}\right)
  = -\frac{1}{4}\cdot 2 F_{\alpha\beta}\frac{\partial F_{\alpha\beta}}{\partial(\partial A_\lambda/\partial x_\mu)}.
\]

Now $\partial F_{\alpha\beta}/\partial(\partial A_\lambda/\partial x_\mu) = \delta_{\alpha\mu}\delta_{\beta\lambda} - \delta_{\alpha\lambda}\delta_{\beta\mu}$. Therefore:
\[
  = -\frac{1}{2}F_{\alpha\beta}(\delta_{\alpha\mu}\delta_{\beta\lambda} - \delta_{\alpha\lambda}\delta_{\beta\mu})
  = -\frac{1}{2}(F_{\mu\lambda} - F_{\lambda\mu})
  = -\frac{1}{2}\cdot 2F_{\mu\lambda} = -F_{\mu\lambda}.
\]
Here we used $F_{\lambda\mu} = -F_{\mu\lambda}$.

\subsubsection*{Step 3: Euler--Lagrange equation}

\[
  \frac{1}{c}j_\lambda - \sum_\mu \frac{\partial}{\partial x_\mu}(-F_{\mu\lambda}) = 0,
\]
\[
  \sum_\mu \frac{\partial F_{\mu\lambda}}{\partial x_\mu} = \frac{1}{c}j_\lambda.
\]

This is the covariant form of the inhomogeneous Maxwell equations:
\[
  \boxed{\sum_\mu \frac{\partial F_{\mu\lambda}}{\partial x_\mu} = \frac{j_\lambda}{c}.}
\]

\subsubsection*{Step 4: Recovery of familiar Maxwell equations}

\noindent\textbf{For $\lambda = 1,2,3$ (spatial components):}

Taking $\lambda = 1$ as an example:
\[
  \frac{\partial F_{21}}{\partial x_2} + \frac{\partial F_{31}}{\partial x_3} + \frac{\partial F_{41}}{\partial x_4} = \frac{j_1}{c}.
\]

Using $F_{21} = \partial A_1/\partial x_2 - \partial A_2/\partial x_1 = -B_3$ (the magnetic field component), and $F_{41} = \partial A_1/\partial x_4 - \partial A_4/\partial x_1$ which involves $E_1$. With $x_4 = ict$ and $A_4 = i\phi$:
\[
  F_{41} = \frac{\partial A_1}{\partial(ict)} - \frac{\partial(i\phi)}{\partial x_1} = \frac{1}{ic}\frac{\partial A_1}{\partial t} - i\frac{\partial\phi}{\partial x_1} = -\frac{i}{c}\left(\frac{\partial A_1}{\partial t} + c\frac{\partial\phi}{\partial x_1}\right) = \frac{i}{c}E_1.
\]

So the spatial equations give:
\[
  (\nabla\times{\bf B})_\lambda - \frac{1}{c}\frac{\partial E_\lambda}{\partial t} = \frac{j_\lambda}{c},
\]
i.e.,
\[
  \nabla\times{\bf B} - \frac{1}{c}\frac{\partial{\bf E}}{\partial t} = \frac{{\bf j}}{c},
\]
which is the Amp\`ere--Maxwell law (in Gaussian units).

\medskip
\noindent\textbf{For $\lambda = 4$:}
\[
  \sum_\mu \frac{\partial F_{\mu 4}}{\partial x_\mu} = \frac{j_4}{c} = \frac{ic\rho}{c} = i\rho.
\]
Since $F_{\mu 4}$ involves $iE_\mu/c$ for the spatial components, this yields:
\[
  \nabla\cdot{\bf E} = \rho,
\]
which is Gauss's law.

\medskip
The homogeneous Maxwell equations ($\nabla\cdot{\bf B}=0$ and Faraday's law) follow identically from the definition $F_{\mu\nu} = \partial A_\nu/\partial x_\mu - \partial A_\mu/\partial x_\nu$ (the Bianchi identity) and do not require the variational principle.

\end{document}
